% Figura: Lowering Example - Da C a Quantum Operations
% Includere con: % Figura: Lowering Example - Da C a Quantum Operations
% Includere con: % Figura: Lowering Example - Da C a Quantum Operations
% Includere con: % Figura: Lowering Example - Da C a Quantum Operations
% Includere con: \input{figures/lowering-example}
% Richiede: \usepackage{tikz}

\begin{figure}[!htb]
\centering
\resizebox{\columnwidth}{!}{%
\begin{tikzpicture}[
    box/.style={
        draw,
        rounded corners=2pt,
        minimum width=2.4cm,
        align=left,
        font=\ttfamily\tiny,
        inner sep=3pt
    },
    header/.style={
        font=\sffamily\scriptsize\bfseries,
        minimum width=2.4cm,
        align=center
    },
    arrow/.style={
        ->,
        >=stealth,
        thick
    }
]

% Colonna 1: C Source
\node[header] (h1) at (0, 0) {C Source};
\node[box, below=0.2cm of h1, text width=2.2cm] (c1) {
int a[2]=\{1,2\};\\
int b[2]=\{3,4\};\\
int c[2];\\
for(i=0;i<2;i++)\\
\ c[i]=a[i]+b[i];
};

% Colonna 2: VecMat IR
\node[header] (h2) at (3, 0) {VecMat IR};
\node[box, below=0.2cm of h2, text width=2.2cm] (c2) {
VecAdd \{\\
\ dest: "c"\\
\ lhs: "a"\\
\ rhs: "b"\\
\ length: 2\\
\}
};

% Colonna 3: QAR
\node[header] (h3) at (6, 0) {QAR};
\node[box, below=0.2cm of h3, text width=2.2cm] (c3) {
QarMapAdd \{\\
\ dest: "c"\\
\ lhs: "a"\\
\ rhs: "b"\\
\ length: 2\\
\}
};

% Colonna 4: Quantum Ops
\node[header] (h4) at (9, 0) {Quantum Ops};
\node[box, below=0.2cm of h4, text width=2.2cm] (c4) {
\%a0 = q.init 1\\
\%a1 = q.init 2\\
\%b0 = q.init 3\\
\%b1 = q.init 4\\
\%c0 = q.init 0\\
\%c1 = q.init 0\\
\%c0' = q.addi \%a0,\%b0\\
\%c1' = q.addi \%a1,\%b1
};

% Frecce
\draw[arrow] (c1.east) -- (c2.west);
\draw[arrow] (c2.east) -- (c3.west);
\draw[arrow] (c3.east) -- (c4.west);

\end{tikzpicture}%
}
\caption{Lowering of a vector addition from C source through VecMat IR and QAR to quantum operations.}
\label{fig:lowering-example}
\end{figure}

% Richiede: \usepackage{tikz}

\begin{figure}[!htb]
\centering
\resizebox{\columnwidth}{!}{%
\begin{tikzpicture}[
    box/.style={
        draw,
        rounded corners=2pt,
        minimum width=2.4cm,
        align=left,
        font=\ttfamily\tiny,
        inner sep=3pt
    },
    header/.style={
        font=\sffamily\scriptsize\bfseries,
        minimum width=2.4cm,
        align=center
    },
    arrow/.style={
        ->,
        >=stealth,
        thick
    }
]

% Colonna 1: C Source
\node[header] (h1) at (0, 0) {C Source};
\node[box, below=0.2cm of h1, text width=2.2cm] (c1) {
int a[2]=\{1,2\};\\
int b[2]=\{3,4\};\\
int c[2];\\
for(i=0;i<2;i++)\\
\ c[i]=a[i]+b[i];
};

% Colonna 2: VecMat IR
\node[header] (h2) at (3, 0) {VecMat IR};
\node[box, below=0.2cm of h2, text width=2.2cm] (c2) {
VecAdd \{\\
\ dest: "c"\\
\ lhs: "a"\\
\ rhs: "b"\\
\ length: 2\\
\}
};

% Colonna 3: QAR
\node[header] (h3) at (6, 0) {QAR};
\node[box, below=0.2cm of h3, text width=2.2cm] (c3) {
QarMapAdd \{\\
\ dest: "c"\\
\ lhs: "a"\\
\ rhs: "b"\\
\ length: 2\\
\}
};

% Colonna 4: Quantum Ops
\node[header] (h4) at (9, 0) {Quantum Ops};
\node[box, below=0.2cm of h4, text width=2.2cm] (c4) {
\%a0 = q.init 1\\
\%a1 = q.init 2\\
\%b0 = q.init 3\\
\%b1 = q.init 4\\
\%c0 = q.init 0\\
\%c1 = q.init 0\\
\%c0' = q.addi \%a0,\%b0\\
\%c1' = q.addi \%a1,\%b1
};

% Frecce
\draw[arrow] (c1.east) -- (c2.west);
\draw[arrow] (c2.east) -- (c3.west);
\draw[arrow] (c3.east) -- (c4.west);

\end{tikzpicture}%
}
\caption{Lowering of a vector addition from C source through VecMat IR and QAR to quantum operations.}
\label{fig:lowering-example}
\end{figure}

% Richiede: \usepackage{tikz}

\begin{figure}[!htb]
\centering
\resizebox{\columnwidth}{!}{%
\begin{tikzpicture}[
    box/.style={
        draw,
        rounded corners=2pt,
        minimum width=2.4cm,
        align=left,
        font=\ttfamily\tiny,
        inner sep=3pt
    },
    header/.style={
        font=\sffamily\scriptsize\bfseries,
        minimum width=2.4cm,
        align=center
    },
    arrow/.style={
        ->,
        >=stealth,
        thick
    }
]

% Colonna 1: C Source
\node[header] (h1) at (0, 0) {C Source};
\node[box, below=0.2cm of h1, text width=2.2cm] (c1) {
int a[2]=\{1,2\};\\
int b[2]=\{3,4\};\\
int c[2];\\
for(i=0;i<2;i++)\\
\ c[i]=a[i]+b[i];
};

% Colonna 2: VecMat IR
\node[header] (h2) at (3, 0) {VecMat IR};
\node[box, below=0.2cm of h2, text width=2.2cm] (c2) {
VecAdd \{\\
\ dest: "c"\\
\ lhs: "a"\\
\ rhs: "b"\\
\ length: 2\\
\}
};

% Colonna 3: QAR
\node[header] (h3) at (6, 0) {QAR};
\node[box, below=0.2cm of h3, text width=2.2cm] (c3) {
QarMapAdd \{\\
\ dest: "c"\\
\ lhs: "a"\\
\ rhs: "b"\\
\ length: 2\\
\}
};

% Colonna 4: Quantum Ops
\node[header] (h4) at (9, 0) {Quantum Ops};
\node[box, below=0.2cm of h4, text width=2.2cm] (c4) {
\%a0 = q.init 1\\
\%a1 = q.init 2\\
\%b0 = q.init 3\\
\%b1 = q.init 4\\
\%c0 = q.init 0\\
\%c1 = q.init 0\\
\%c0' = q.addi \%a0,\%b0\\
\%c1' = q.addi \%a1,\%b1
};

% Frecce
\draw[arrow] (c1.east) -- (c2.west);
\draw[arrow] (c2.east) -- (c3.west);
\draw[arrow] (c3.east) -- (c4.west);

\end{tikzpicture}%
}
\caption{Lowering of a vector addition from C source through VecMat IR and QAR to quantum operations.}
\label{fig:lowering-example}
\end{figure}

% Richiede: \usepackage{tikz}

\begin{figure}[!htb]
\centering
\resizebox{\columnwidth}{!}{%
\begin{tikzpicture}[
    box/.style={
        draw,
        rounded corners=2pt,
        minimum width=2.4cm,
        align=left,
        font=\ttfamily\tiny,
        inner sep=3pt
    },
    header/.style={
        font=\sffamily\scriptsize\bfseries,
        minimum width=2.4cm,
        align=center
    },
    arrow/.style={
        ->,
        >=stealth,
        thick
    }
]

% Colonna 1: C Source
\node[header] (h1) at (0, 0) {C Source};
\node[box, below=0.2cm of h1, text width=2.2cm] (c1) {
int a[2]=\{1,2\};\\
int b[2]=\{3,4\};\\
int c[2];\\
for(i=0;i<2;i++)\\
\ c[i]=a[i]+b[i];
};

% Colonna 2: VecMat IR
\node[header] (h2) at (3, 0) {VecMat IR};
\node[box, below=0.2cm of h2, text width=2.2cm] (c2) {
VecAdd \{\\
\ dest: "c"\\
\ lhs: "a"\\
\ rhs: "b"\\
\ length: 2\\
\}
};

% Colonna 3: QAR
\node[header] (h3) at (6, 0) {QAR};
\node[box, below=0.2cm of h3, text width=2.2cm] (c3) {
QarMapAdd \{\\
\ dest: "c"\\
\ lhs: "a"\\
\ rhs: "b"\\
\ length: 2\\
\}
};

% Colonna 4: Quantum Ops
\node[header] (h4) at (9, 0) {Quantum Ops};
\node[box, below=0.2cm of h4, text width=2.2cm] (c4) {
\%a0 = q.init 1\\
\%a1 = q.init 2\\
\%b0 = q.init 3\\
\%b1 = q.init 4\\
\%c0 = q.init 0\\
\%c1 = q.init 0\\
\%c0' = q.addi \%a0,\%b0\\
\%c1' = q.addi \%a1,\%b1
};

% Frecce
\draw[arrow] (c1.east) -- (c2.west);
\draw[arrow] (c2.east) -- (c3.west);
\draw[arrow] (c3.east) -- (c4.west);

\end{tikzpicture}%
}
\caption{Lowering of a vector addition from C source through VecMat IR and QAR to quantum operations.}
\label{fig:lowering-example}
\end{figure}
